\documentclass[10pt,a4paper]{amsart}
\usepackage{amsmath,amssymb,amsthm,graphicx}
% Retained preamble for the Stacks Project chapter snapshot in this directory.
% The chapter's own preamble.tex is a copy of the shared Stacks Project preamble; the block
% below is sliced MECHANICALLY, line for line, from it. Nothing was added, removed or
% redefined; the wrapper class, the xr-hyper/hyperref wiring and the externaldocument list
% are omitted because the assembler supplies the document class.
%
% lines taken from preamble.tex: 3, 5, 9, 11, 12, 156, 157, 158, ... (51 lines total)
%
% package coverage restored from the chapter's own preamble (2026-10-13): xy, xy-2cell

\usepackage[all]{xy}
\xyoption{2cell}
\UseAllTwocells

\newenvironment{reference}{\comment}{\endcomment}
\newenvironment{slogan}{\comment}{\endcomment}
\newenvironment{history}{\comment}{\endcomment}
\theoremstyle{plain}
\newtheorem{theorem}[subsection]{Theorem}
\newtheorem{proposition}[subsection]{Proposition}
\newtheorem{lemma}[subsection]{Lemma}
\theoremstyle{definition}
\newtheorem{definition}[subsection]{Definition}
\newtheorem{example}[subsection]{Example}
\newtheorem{exercise}[subsection]{Exercise}
\newtheorem{situation}[subsection]{Situation}
\theoremstyle{remark}
\newtheorem{remark}[subsection]{Remark}
\newtheorem{remarks}[subsection]{Remarks}
\numberwithin{equation}{subsection}
\def\lim{\mathop{\mathrm{lim}}\nolimits}
\def\colim{\mathop{\mathrm{colim}}\nolimits}
\def\Spec{\mathop{\mathrm{Spec}}}
\def\Hom{\mathop{\mathrm{Hom}}\nolimits}
\def\Ext{\mathop{\mathrm{Ext}}\nolimits}
\def\SheafHom{\mathop{\mathcal{H}\!\mathit{om}}\nolimits}
\def\SheafExt{\mathop{\mathcal{E}\!\mathit{xt}}\nolimits}
\def\Sch{\mathit{Sch}}
\def\Mor{\mathop{\mathrm{Mor}}\nolimits}
\def\Ob{\mathop{\mathrm{Ob}}\nolimits}
\def\Sh{\mathop{\mathit{Sh}}\nolimits}
\def\NL{\mathop{N\!L}\nolimits}
\def\CH{\mathop{\mathrm{CH}}\nolimits}
\def\proetale{{pro\text{-}\acute{e}tale}}
\def\etale{{\acute{e}tale}}
\def\QCoh{\mathit{QCoh}}
\def\Ker{\mathop{\mathrm{Ker}}}
\def\Im{\mathop{\mathrm{Im}}}
\def\Coker{\mathop{\mathrm{Coker}}}
\def\Coim{\mathop{\mathrm{Coim}}}
\DeclareMathSymbol{\boxtimes}{\mathbin}{AMSa}{"02}
\def\QCohstack{\mathcal{QC}\!\mathit{oh}}
\def\Cohstack{\mathcal{C}\!\mathit{oh}}
\def\Spacesstack{\mathcal{S}\!\mathit{paces}}
\def\Quotfunctor{\mathrm{Quot}}
\def\Hilbfunctor{\mathrm{Hilb}}
\def\Curvesstack{\mathcal{C}\!\mathit{urves}}
\def\Polarizedstack{\mathcal{P}\!\mathit{olarized}}
\def\Complexesstack{\mathcal{C}\!\mathit{omplexes}}
\def\Pic{\mathop{\mathrm{Pic}}\nolimits}
\def\Picardstack{\mathcal{P}\!\mathit{ic}}
\def\Picardfunctor{\mathrm{Pic}}
\def\Deformationcategory{\mathcal{D}\!\mathit{ef}}
\usepackage{cleveref}
\crefname{originallemma}{original lemma}{original lemmas}
\Crefname{originallemma}{Original lemma}{Original lemmas}
\makeatletter
\newlabel{bootstrap-lemma-descend-algebraic-space}{{0ADV}{source}{}{sourcepointer-1}{}}
\newlabel{bootstrap-lemma-descend-algebraic-space@cref}{{[originallemma][0][]0ADV}{[1][0][]source}{}{}{}}
\makeatother
\crefname{originallemma}{original lemma}{original lemmas}
\Crefname{originallemma}{Original lemma}{Original lemmas}
\makeatletter
\newlabel{quot-situation-curves}{{0D4Z}{source}{}{sourcepointer-2}{}}
\newlabel{quot-situation-curves@cref}{{[originallemma][0][]0D4Z}{[1][0][]source}{}{}{}}
\makeatother
\crefname{originallemma}{original lemma}{original lemmas}
\Crefname{originallemma}{Original lemma}{Original lemmas}
\makeatletter
\newlabel{spaces-descent-lemma-descending-property-finite-presentation}{{041V}{source}{}{sourcepointer-3}{}}
\newlabel{spaces-descent-lemma-descending-property-finite-presentation@cref}{{[originallemma][0][]041V}{[1][0][]source}{}{}{}}
\makeatother
\crefname{originallemma}{original lemma}{original lemmas}
\Crefname{originallemma}{Original lemma}{Original lemmas}
\makeatletter
\newlabel{spaces-descent-lemma-descending-property-flat}{{041W}{source}{}{sourcepointer-4}{}}
\newlabel{spaces-descent-lemma-descending-property-flat@cref}{{[originallemma][0][]041W}{[1][0][]source}{}{}{}}
\makeatother
\crefname{originallemma}{original lemma}{original lemmas}
\Crefname{originallemma}{Original lemma}{Original lemmas}
\makeatletter
\newlabel{spaces-descent-lemma-descending-property-proper}{{0422}{source}{}{sourcepointer-5}{}}
\newlabel{spaces-descent-lemma-descending-property-proper@cref}{{[originallemma][0][]0422}{[1][0][]source}{}{}{}}
\makeatother
\crefname{originallemma}{original lemma}{original lemmas}
\Crefname{originallemma}{Original lemma}{Original lemmas}
\makeatletter
\newlabel{spaces-groupoids-lemma-quotient-stack-2-coequalizer}{{044U}{source}{}{sourcepointer-6}{}}
\newlabel{spaces-groupoids-lemma-quotient-stack-2-coequalizer@cref}{{[originallemma][0][]044U}{[1][0][]source}{}{}{}}
\makeatother
\crefname{originallemma}{original lemma}{original lemmas}
\Crefname{originallemma}{Original lemma}{Original lemmas}
\makeatletter
\newlabel{spaces-morphisms-definition-relative-dimension}{{06LR}{source}{}{sourcepointer-7}{}}
\newlabel{spaces-morphisms-definition-relative-dimension@cref}{{[originallemma][0][]06LR}{[1][0][]source}{}{}{}}
\makeatother
\title{Stacks, Moduli of Curves: lemma-extend-curves-to-spaces}
\author{The Stacks Project Authors}
\date{stacks-project:moduli-curves.tex:local-snapshot-20261012; GNU FDL 1.2 or later}
\begin{document}
\maketitle
\noindent\textit{Editorial note.} Statement span is lines 298-305 of sources/stacks-moduli-curves/pinned-original-source.tex; proof span is lines 307-362, the complete author proof (54 body lines).\par\medskip
\noindent\textit{Editorial note.} Navigation stubs: 7 cited-but-not-included labels are supplied as non-load-bearing pointers with their pinned tags where the snapshot tag index has them.\par\medskip
\noindent\textit{Editorial note.} extra\_packages is empty on purpose: the retained preamble already loads xy with options and a second declaration clashes in the assembler wrapper.\par\medskip
\noindent\textit{Editorial note.} Text is retained verbatim, including the author's own typos and placeholders.\par\medskip
\noindent\textit{Editorial note.} Queue index 5 of queue\_moduli-curves.json (pb=2671, stmt 298-305, proof 307-362).\par\medskip
\noindent\textit{Allowed standard prerequisites (editorial).}\begin{itemize}
\item Explicit prerequisite boundary. The cited inputs are prerequisites; the argument in the retained span is not replaced by any cited result.
\end{itemize}\medskip
\ifdefined\hypertarget\hypertarget{sourcepointer-1}{}\fi
\noindent\textit{Source reference.} Original lemma 0ADV: sources/stacks-moduli-curves/pinned-original-source.tex, source lines 362--362. Cited by the retained statement or proof but not part of the collected spans; supplied only as a navigation pointer. The argument does not depend on its content.\par\medskip
\ifdefined\hypertarget\hypertarget{sourcepointer-2}{}\fi
\noindent\textit{Source reference.} Original lemma 0D4Z: sources/stacks-moduli-curves/pinned-original-source.tex, source lines 362--362. Cited by the retained statement or proof but not part of the collected spans; supplied only as a navigation pointer. The argument does not depend on its content.\par\medskip
\ifdefined\hypertarget\hypertarget{sourcepointer-3}{}\fi
\noindent\textit{Source reference.} Original lemma 041V: sources/stacks-moduli-curves/pinned-original-source.tex, source lines 362--362. Cited by the retained statement or proof but not part of the collected spans; supplied only as a navigation pointer. The argument does not depend on its content.\par\medskip
\ifdefined\hypertarget\hypertarget{sourcepointer-4}{}\fi
\noindent\textit{Source reference.} Original lemma 041W: sources/stacks-moduli-curves/pinned-original-source.tex, source lines 362--362. Cited by the retained statement or proof but not part of the collected spans; supplied only as a navigation pointer. The argument does not depend on its content.\par\medskip
\ifdefined\hypertarget\hypertarget{sourcepointer-5}{}\fi
\noindent\textit{Source reference.} Original lemma 0422: sources/stacks-moduli-curves/pinned-original-source.tex, source lines 362--362. Cited by the retained statement or proof but not part of the collected spans; supplied only as a navigation pointer. The argument does not depend on its content.\par\medskip
\ifdefined\hypertarget\hypertarget{sourcepointer-6}{}\fi
\noindent\textit{Source reference.} Original lemma 044U: sources/stacks-moduli-curves/pinned-original-source.tex, source lines 362--362. Cited by the retained statement or proof but not part of the collected spans; supplied only as a navigation pointer. The argument does not depend on its content.\par\medskip
\ifdefined\hypertarget\hypertarget{sourcepointer-7}{}\fi
\noindent\textit{Source reference.} Original lemma 06LR: sources/stacks-moduli-curves/pinned-original-source.tex, source lines 362--362. Cited by the retained statement or proof but not part of the collected spans; supplied only as a navigation pointer. The argument does not depend on its content.\par\medskip
\begin{lemma}
\label{lemma-extend-curves-to-spaces}
Let $T \to B$ be a morphism of algebraic spaces. The category
$$
\Mor_B(T, B\text{-}\Curvesstack) = \Mor(T, \Curvesstack)
$$
is the category of families of curves over $T$.
\end{lemma}

\begin{proof}
A family of curves over $T$ is a morphism $f : X \to T$ of algebraic
spaces, which is flat, proper, of finite presentation, and has
relative dimension $\leq 1$ (Morphisms of Spaces, Definition
\ref{spaces-morphisms-definition-relative-dimension}).
This is exactly the same as the definition in
Quot, Situation \ref{quot-situation-curves}
except that $T$ the base is allowed to be an algebraic space.
Our default base category for algebraic stacks/spaces
is the category of schemes, hence the lemma does not follow
immediately from the definitions. Having said this, we encourage
the reader to skip the proof.

\medskip\noindent
By the product description of $B\text{-}\Curvesstack$ given above,
it suffices to prove the lemma in the absolute case. Choose a scheme
$U$ and a surjective \'etale morphism $p : U \to T$.
Let $R = U \times_T U$ with projections $s, t : R \to U$.

\medskip\noindent
Let $v : T \to \Curvesstack$ be a morphism. Then $v \circ p$
corresponds to a family of curves $X_U \to U$. The canonical
$2$-morphism $v \circ p \circ t \to v \circ p \circ s$
is an isomorphism $\varphi : X_U \times_{U, s} R \to X_U \times_{U, t} R$.
This isomorphism satisfies the cocycle condition on
$R \times_{s, t} R$.
By Bootstrap, Lemma \ref{bootstrap-lemma-descend-algebraic-space}
we obtain a morphism of algebraic spaces $X \to T$
whose pullback to $U$ is equal to $X_U$ compatible with $\varphi$.
Since $\{U \to T\}$ is an \'etale covering, we see that
$X \to T$ is flat, proper, of finite presentation by
Descent on Spaces, Lemmas
\ref{spaces-descent-lemma-descending-property-flat},
\ref{spaces-descent-lemma-descending-property-proper}, and
\ref{spaces-descent-lemma-descending-property-finite-presentation}.
Also $X \to T$ has relative dimension $\leq 1$ because this is
an \'etale local property. Hence $X \to T$ is a family of curves over $T$.

\medskip\noindent
Conversely, let $X \to T$ be a family of curves. Then the
base change $X_U$ determines a morphism $w : U \to \Curvesstack$
and the canonical isomorphism $X_U \times_{U, s} R \to X_U \times_{U, t} R$
determines a $2$-arrow $w \circ s \to w \circ t$ satisfying the
cocycle condition. Thus a morphism $v : T = [U/R] \to \Curvesstack$
by the universal property of the quotient $[U/R]$, see
Groupoids in Spaces, Lemma
\ref{spaces-groupoids-lemma-quotient-stack-2-coequalizer}.
(Actually, it is much easier in this case to go back to before
we introduced our abuse of language and direct construct
the functor $\Sch/T \to \Curvesstack$ which ``is'' the
morphism $T \to \Curvesstack$.)

\medskip\noindent
We omit the verification that the constructions given above
extend to morphisms between objects and are mutually quasi-inverse.
\end{proof}

\end{document}
